\frontchapter{Abstract}

Automated ball-launching machines are widely used in sports training, but most practical systems operate in open-loop mode: the trainer configures a fixed angle, speed, and interval, and the machine repeats that program without measuring the athlete's current body position. This thesis presents the design, implementation, and partial validation of a vision-guided ball-launching system that estimates selected human body joints (right knee, right hip, and left shoulder) in three dimensions and computes corresponding aim parameters for a Ball Launching Machine (BLM) in a domestic garage arena. The current work validates the perception, calibration, targeting, and safety-gated integration components, but it does not claim a completed full closed-loop moving-target firing demonstration.

The system uses four fixed commodity USB cameras calibrated with ChArUco boards for intrinsic parameters and AprilTag fiducial markers for extrinsic world-frame registration. Ball detection is performed with a YOLO-based detector, and human pose estimation is performed with an open-source pose-estimation backend using the COCO 17-keypoint skeleton convention. Multi-view triangulation converts synchronized two-dimensional observations into three-dimensional world-frame coordinates in millimetres. A ballistic solver computes pitch, yaw, and wheel-speed commands from the selected target joint, while low-level actuation is handled by an ESP32-based controller. A staged safety protocol, including E-STOP handling, command gating, and structured decision logging, is used to restrict actuation during integration.

Quantitative evaluation is based on a 36-point static ball ground-truth dataset and an 81-trial joint-touch dataset. Under the current repository calibration bundle, the raw static ball evaluation gives a mean 3D error of 156.90 mm and a P95 error of 288.34 mm, with high temporal precision (3.09 mm mean standard-deviation norm). The joint-touch evaluation gives a mean 3D error of 178.98 mm and a P95 error of 243.77 mm over 62 valid trials. These results show that the system provides a repeatable and practically useful perception backbone, while also revealing systematic calibration bias, occlusion sensitivity, and the need for further closed-loop validation. The thesis therefore positions the work as a strong but incomplete step toward affordable pose-reactive ball delivery, with moving-target firing and independent validation left as future work.
